\section*{摘要 Abstract}

现代恒星巡天测量着数以百万计的光谱，然而计算一个自洽的大气模型及其光谱可能需要数十分钟。这一计算成本催生了预计算网格、光谱模拟器和数据驱动模型。我们提出 \paynezero：它将一维 LTE Kurucz 计算重组为 GPU 原生的光谱合成与多核 CPU 大气迭代，并对照原始 Fortran 程序完成验证。在 NVIDIA H100 GPU 上，以 $R_{\rm grid}=300{,}000$ 采样的 300--1000 nm 太阳光谱约需 14 秒；APOGEE 1500--1700 nm 区间约需 1 秒。物理大气迭代在 16 线程 AMD CPU 上需 2--5 秒，学习型初始化器进一步减少了收敛所需的迭代次数。在所测试的矮星与巨星参数区间内，最终光谱与原版计算保持实用层面的一致。这样的速度使直接合成可以被放进优化器内部，而无需“标签到流量”的光谱模拟器。我们演示了对约减后 APOGEE 光谱的多元素直接拟合，恢复出与巡天星表大体一致的多元素丰度趋势。GPU 常驻的视向速度平移、致宽、线扩展函数卷积与探测器采样相对合成的额外开销可以忽略。直接合成搜索在 H100 上每颗星耗时不到一分钟，而大气验证在多核 CPU 上独立运行。同一计算图还能针对太阳和大角星联合定标超过 $10^5$ 个振子强度与阻尼改正，在 H100 上约需一分钟。因此，\paynezero\ 把直接物理拟合与原子数据定标推进到了巡天规模。代码开源于 \url{https://github.com/tingyuansen/payne-zero}。

\begin{EN}
Modern stellar surveys measure millions of spectra, yet one self-consistent
atmosphere and spectrum can require tens of minutes. This cost has motivated
grids, spectral emulators, and data-driven models. We present \paynezero, which
reorganizes one-dimensional LTE Kurucz calculations for GPU-native synthesis
and multicore atmosphere iteration, and validate it against the original
Fortran programs. A
300--1000~nm solar spectrum sampled at $R_{\rm grid}=300{,}000$ takes about
14~s on an NVIDIA H100 GPU, while the APOGEE 1500--1700~nm interval takes about
1~s. Physical atmosphere iterations take 2--5~s on 16 AMD CPU threads, and
learned initializers reduce the iterations required for convergence. Final
spectra remain in practical parity across the tested dwarf and giant regimes.
These speeds place direct synthesis inside an optimizer without a
label-to-flux spectral emulator. We demonstrate direct many-element fitting of
reduced APOGEE spectra and recover multi-element abundance trends broadly
consistent with the survey catalog.
GPU-resident velocity shifts, broadening, line-spread-function convolution,
and detector sampling add negligible cost relative to synthesis.
The direct-synthesis search takes less than one minute per star on an H100, while
atmosphere verification runs independently on multicore CPUs. The same
computational graph calibrates more than $10^5$ oscillator-strength and
damping corrections jointly against the Sun and
Arcturus in about one minute on an H100. \paynezero\ therefore brings direct
physical fitting and atomic-data calibration to survey scale. The code is
available at \url{https://github.com/tingyuansen/payne-zero}.
\end{EN}

\noindent\textbf{关键词（原文）：}恒星大气；恒星光谱线；天文软件；天文数据分析；计算方法。

